% GENERATED FILE. Edit canonical lessons, then run npm run generate:books.
% canonical-source-hash: fnv1a64:e6496179200cc39f
% canonical-lessons: RU-C318-traditsiya, RU-C318-obychay, RU-C318-obryad, RU-C318-bal, RU-C318-gulyanye, RU-R318-review-311-318

\chapter{Customs and festivities}
\label{ch:ru-customs-and-festivities}
\hlchaptermodality{318}

\begin{chapteropening}
\textbf{By the end of this chapter:} \emph{I can talk about a tradition, a custom, a rite, a ball and public festivities.}

Complete the last lesson of chapter 318: I can talk about a tradition, a custom, a rite, a ball and public festivities.
\end{chapteropening}

\section[\texorpdfstring{\ru{традиция} (tradítsiya)}{традиция (tradítsiya)}]{\ru{традиция} (tradítsiya) — a tradition}

\label{lesson:RU-C318-traditsiya}



\begin{quote}
Before the new one: say the Russian for to blow out, then the Russian for to have fun.
\end{quote}

\subsection*{You'll want to know: \ru{традиция}}
\textbf{\ru{традиция}} (\emph{tradítsiya}) — \textquotedblleft{}a tradition\textquotedblright{}.

\begin{quote}
\ru{В} \ru{нашей} \ru{семье} \ru{есть} \ru{старая} \ru{традиция}.
\end{quote}

\begin{quote}
\emph{V náshey semyé yest stáraya tradítsiya.}
\end{quote}

\begin{quote}
Our family has an old tradition.
\end{quote}

Feminine. \emph{\ru{По} \ru{традиции}} — by tradition.

\begin{cousinweb}
\emph{\ru{традиция}} goes back to Latin \emph{trāditiō} \textquotedblleft{}a handing over\textquotedblright{} (EtymDB), from \emph{trādere} \textquotedblleft{}to hand over\textquotedblright{} (Lewis \& Short).
\end{cousinweb}

\begin{grammarlens}[title={What you've built}]
One more word for telling, explaining and describing at B1.
\end{grammarlens}

\subsection*{Your turn}
\begin{itemize}
\raggedright
  \item \emph{Say it:} \emph{tradítsiya}
  \item \emph{Say it:} \emph{tradítsiya}, once more
  \item \emph{Say it:} a sentence of your own with \emph{tradítsiya}
  \item \emph{Recall:} say the Russian for to blow out, then the Russian for to have fun, then say \emph{tradítsiya} again
\end{itemize}

\begin{tcolorbox}[breakable,colback=teal!4,colframe=teal!35!black,title={Before you move on}]
What does \ru{традиция} mean? (\textquotedblleft{}A tradition\textquotedblright{}.) Say it once more.
\end{tcolorbox}
\section[\texorpdfstring{\ru{обычай} (obýchay)}{обычай (obýchay)}]{\ru{обычай} (obýchay) — a custom}

\label{lesson:RU-C318-obychay}



\begin{quote}
Before the new one: say the Russian for to have fun, then the Russian for a tradition.
\end{quote}

\subsection*{You'll want to know: \ru{обычай}}
\textbf{\ru{обычай}} (\emph{obýchay}) — \textquotedblleft{}a custom\textquotedblright{}.

\begin{quote}
\ru{По} \ru{русскому} \ru{обычаю} \ru{гостей} \ru{встречают} \ru{хлебом} \ru{и} \ru{солью}.
\end{quote}

\begin{quote}
\emph{Po rússkomu obýchayu gostéy vstrecháyut khlébom i sólyu.}
\end{quote}

\begin{quote}
By Russian custom, guests are welcomed with bread and salt.
\end{quote}

Masculine.

\begin{cousinweb}
\emph{\ru{обычай}} is inherited from Proto-Slavic \emph{obvyčajĭ} (EtymDB), built on \emph{vyknǫti} \textquotedblleft{}to get used to\textquotedblright{}, the root of \emph{\ru{привыкать}}.
\end{cousinweb}

\begin{grammarlens}[title={What you've built}]
One more word for telling, explaining and describing at B1.
\end{grammarlens}

\subsection*{Your turn}
\begin{itemize}
\raggedright
  \item \emph{Say it:} \emph{obýchay}
  \item \emph{Say it:} \emph{obýchay}, once more
  \item \emph{Say it:} a sentence of your own with \emph{obýchay}
  \item \emph{Recall:} say the Russian for to have fun, then the Russian for a tradition, then say \emph{obýchay} again
\end{itemize}

\begin{tcolorbox}[breakable,colback=teal!4,colframe=teal!35!black,title={Before you move on}]
What does \ru{обычай} mean? (\textquotedblleft{}A custom\textquotedblright{}.) Say it once more.
\end{tcolorbox}
\section[\texorpdfstring{\ru{обряд} (obryád)}{обряд (obryád)}]{\ru{обряд} (obryád) — a rite, a ritual}

\label{lesson:RU-C318-obryad}



\begin{quote}
Before the new one: say the Russian for a tradition, then the Russian for a custom.
\end{quote}

\subsection*{You'll want to know: \ru{обряд}}
\textbf{\ru{обряд}} (\emph{obryád}) — \textquotedblleft{}a rite, a ritual\textquotedblright{}.

\begin{quote}
\ru{Это} \ru{очень} \ru{старый} \ru{обряд}.
\end{quote}

\begin{quote}
\emph{Éto óchen stáryy obryád.}
\end{quote}

\begin{quote}
This is a very old rite.
\end{quote}

Masculine.

\begin{cousinweb}
\emph{\ru{обряд}} is usually traced to \emph{\ru{об}-} and \emph{\ru{ряд}} \textquotedblleft{}a row, order\textquotedblright{}: a set order of acts. EtymDB lists no parent.
\end{cousinweb}

\begin{grammarlens}[title={What you've built}]
One more word for telling, explaining and describing at B1.
\end{grammarlens}

\subsection*{Your turn}
\begin{itemize}
\raggedright
  \item \emph{Say it:} \emph{obryád}
  \item \emph{Say it:} \emph{obryád}, once more
  \item \emph{Say it:} a sentence of your own with \emph{obryád}
  \item \emph{Recall:} say the Russian for a tradition, then the Russian for a custom, then say \emph{obryád} again
\end{itemize}

\begin{tcolorbox}[breakable,colback=teal!4,colframe=teal!35!black,title={Before you move on}]
What does \ru{обряд} mean? (\textquotedblleft{}A rite, a ritual\textquotedblright{}.) Say it once more.
\end{tcolorbox}
\section[\texorpdfstring{\ru{бал} (bal)}{бал (bal)}]{\ru{бал} (bal) — a ball (a dance)}

\label{lesson:RU-C318-bal}



\begin{quote}
Before the new one: say the Russian for a custom, then the Russian for a rite.
\end{quote}

\subsection*{You'll want to know: \ru{бал}}
\textbf{\ru{бал}} (\emph{bal}) — \textquotedblleft{}a ball (a dance)\textquotedblright{}.

\begin{quote}
\ru{Студенты} \ru{готовились} \ru{к} \ru{новогоднему} \ru{балу}.
\end{quote}

\begin{quote}
\emph{Studénty gotóvilis k novogódnemu bálu.}
\end{quote}

\begin{quote}
The students were getting ready for the New Year ball.
\end{quote}

\emph{\ru{На} \ru{балу}} — at the ball.

\begin{cousinweb}
\emph{\ru{бал}} is usually traced through German \emph{Ball} or French \emph{bal} to Late Latin \emph{ballāre} \textquotedblleft{}to dance\textquotedblright{}; EtymDB lists no parent.
\end{cousinweb}

\begin{grammarlens}[title={What you've built}]
One more word for telling, explaining and describing at B1.
\end{grammarlens}

\subsection*{Your turn}
\begin{itemize}
\raggedright
  \item \emph{Say it:} \emph{bal}
  \item \emph{Say it:} \emph{bal}, once more
  \item \emph{Say it:} a sentence of your own with \emph{bal}
  \item \emph{Recall:} say the Russian for a custom, then the Russian for a rite, then say \emph{bal} again
\end{itemize}

\begin{tcolorbox}[breakable,colback=teal!4,colframe=teal!35!black,title={Before you move on}]
What does \ru{бал} mean? (\textquotedblleft{}A ball (a dance)\textquotedblright{}.) Say it once more.
\end{tcolorbox}
\section[\texorpdfstring{\ru{гулянье} (gulyánye)}{гулянье (gulyánye)}]{\ru{гулянье} (gulyánye) — festivities, a fête}

\label{lesson:RU-C318-gulyanye}



\begin{quote}
Before the new one: say the Russian for a rite, then the Russian for a ball.
\end{quote}

\subsection*{You'll want to know: \ru{гулянье}}
\textbf{\ru{гулянье}} (\emph{gulyánye}) — \textquotedblleft{}festivities, a fête\textquotedblright{}.

\begin{quote}
\ru{На} \ru{площади} \ru{было} \ru{большое} \ru{гулянье}.
\end{quote}

\begin{quote}
\emph{Na plóshchadi býlo bolshóye gulyánye.}
\end{quote}

\begin{quote}
There were big festivities on the square.
\end{quote}

Neuter; also spelt \emph{\ru{гуляние}}. \emph{\ru{Народное} \ru{гулянье}} — public festivities.

\begin{cousinweb}
\emph{\ru{гулянье}} is built on \emph{\ru{гулять}} \textquotedblleft{}to stroll, to make merry\textquotedblright{}, an East Slavic verb whose origin is not settled; EtymDB lists no parent.
\end{cousinweb}

\begin{grammarlens}[title={What you've built}]
One more word for telling, explaining and describing at B1.
\end{grammarlens}

\subsection*{Your turn}
\begin{itemize}
\raggedright
  \item \emph{Say it:} \emph{gulyánye}
  \item \emph{Say it:} \emph{gulyánye}, once more
  \item \emph{Say it:} a sentence of your own with \emph{gulyánye}
  \item \emph{Recall:} say the Russian for a rite, then the Russian for a ball, then say \emph{gulyánye} again
\end{itemize}

\begin{tcolorbox}[breakable,colback=teal!4,colframe=teal!35!black,title={Before you move on}]
What does \ru{гулянье} mean? (\textquotedblleft{}Festivities, a fête\textquotedblright{}.) Say it once more.
\end{tcolorbox}
\section[(dialogue)]{Words from chapters 311–318 — a review}

\label{lesson:RU-R318-review-311-318}



\begin{quote}
Say the Russian for to celebrate and to mark.
\end{quote}

\subsection*{Your turn}
Say these aloud:

\begin{itemize}
\raggedright
  \item to celebrate, to mark, a celebration, solemn, festive
  \item to congratulate, congratulations, a wish, to wish, a toast
  \item a New Year tree, fairy lights, to decorate, a decoration, to dress up
  \item a firework display, fireworks, confetti, a party popper, a balloon
  \item treats, to treat, a feast, to clink glasses, champagne
  \item a carnival, a masquerade, a parade, a procession, a fair
  \item the birthday boy, a name day, to make, to blow out, to have fun
  \item a tradition, a custom, a rite, a ball, festivities
\end{itemize}

\begin{tcolorbox}[breakable,colback=teal!4,colframe=teal!35!black,title={Before you move on}]
To celebrate? (\textbf{\ru{праздновать}}.) To mark? (\textbf{\ru{отмечать}}.) A celebration? (\textbf{\ru{торжество}}.) Solemn? (\textbf{\ru{торжественный}}.) Festive? (\textbf{\ru{праздничный}}.) To congratulate? (\textbf{\ru{поздравлять}}.) Congratulations? (\textbf{\ru{поздравление}}.) A wish? (\textbf{\ru{пожелание}}.) To wish? (\textbf{\ru{желать}}.) A toast? (\textbf{\ru{тост}}.) A New Year tree? (\textbf{\ru{ёлка}}.) Fairy lights? (\textbf{\ru{гирлянда}}.) To decorate? (\textbf{\ru{украшать}}.) A decoration? (\textbf{\ru{украшение}}.) To dress up? (\textbf{\ru{наряжать}}.) A firework display? (\textbf{\ru{салют}}.) Fireworks? (\textbf{\ru{фейерверк}}.) Confetti? (\textbf{\ru{конфетти}}.) A party popper? (\textbf{\ru{хлопушка}}.) A balloon? (\textbf{\ru{шарик}}.) Treats? (\textbf{\ru{угощение}}.) To treat? (\textbf{\ru{угощать}}.) A feast? (\textbf{\ru{застолье}}.) To clink glasses? (\textbf{\ru{чокаться}}.) Champagne? (\textbf{\ru{шампанское}}.) A carnival? (\textbf{\ru{карнавал}}.) A masquerade? (\textbf{\ru{маскарад}}.) A parade? (\textbf{\ru{парад}}.) A procession? (\textbf{\ru{шествие}}.) A fair? (\textbf{\ru{ярмарка}}.) The birthday boy? (\textbf{\ru{именинник}}.) A name day? (\textbf{\ru{именины}}.) To make? (\textbf{\ru{загадывать}}.) To blow out? (\textbf{\ru{задувать}}.) To have fun? (\textbf{\ru{веселиться}}.) A tradition? (\textbf{\ru{традиция}}.) A custom? (\textbf{\ru{обычай}}.) A rite? (\textbf{\ru{обряд}}.) A ball? (\textbf{\ru{бал}}.) Festivities? (\textbf{\ru{гулянье}}.)
\end{tcolorbox}
